\section{Methods (Prespecified)}
This pre-analysis protocol is reported following STROBE for the observational analyses \citep{vonelm2007strobe} and TRIPOD+AI for the prediction models \citep{collins2024tripodai}; the prospective clinical extension will follow SPIRIT. A completed reporting checklist accompanies the repository. All design choices below are fixed prior to any real-data analysis.

\subsection{Study Design}
The work has three stages, summarized in the conceptual framework of Figure~\ref{fig:framework}. \textbf{Stage~1} analyzes KNHANES as the primary public-data development cohort (association estimation and research-model development). \textbf{Stage~2} evaluates cross-national transportability in NHANES (geographic external validation). \textbf{Stage~3} prepares a prospective clinical extension at Ajou University Hospital for clinical external validation, treatment-timing analysis in SSNHL, and evaluation of the AI extraction and safety components. Stages~1--2 use only deidentified public-use files; Stage~3 proceeds only under IRB approval.

\subsection{Fixed Datasets, Cycles, and Eligibility}
To eliminate analyst degrees of freedom, the survey cycles, files, ages, frequencies, and exclusions are fixed in advance (Table~\ref{tab:datasets}). \textbf{KNHANES:} the otologic-examination cycles carrying pure-tone audiometry and tinnitus items (KNHANES IV--V, 2009--2012) form the primary development window; cycles are pooled only after confirming identical item wording and audiometry protocols. \textbf{NHANES:} the audiometry-module cycles 2011--2012, 2015--2016, and 2017--March~2020 form the external-validation window. For both, eligible records are adults aged 40--69~years (the band jointly covered by the audiometry modules and our target of mid-to-late-career workers), with valid examination weights and a non-missing value for the outcome under analysis. Pure-tone thresholds are analyzed at 0.5, 1, 2, 3, 4, 6, and 8~kHz per ear; conductive patterns are flagged and, where tympanometry/air--bone data permit, excluded in a sensitivity analysis. We prespecify expected raw and analytic sample sizes per window (reported in the repository data dictionary) and the exact set of predictors available in \emph{both} surveys; any predictor available in only one survey is excluded from the common transportable model and used only in country-specific extended models.

\subsection{Exposure: Perceived Stress and Work-Related Factors}
Consistent with the naming principle (Introduction), the primary exposure is \emph{perceived stress}, operationalized from each survey's general stress/psychological-distress item, and \emph{work-related factors} (occupation category, employment status, and---where a cycle provides them---working hours and shift work) are modeled as separate covariates or effect modifiers, never as a proxy for validated job strain. Because item wording, response categories, and cultural response styles differ across surveys, Table~\ref{tab:harmonization} records, for every exposure and covariate, the original question, response scale, harmonization rule, and a measurement-comparability rating (high/medium/low). Variables rated \emph{low} comparability (typically the stress and some psychosocial items) are (a) excluded from the common transportable model, (b) entered only in country-specific models, and (c) examined in a formal measurement-invariance check (configural/metric where a latent construct is defined). We do not pool a variable across countries unless at least metric-level comparability is plausible.

\subsection{Causal Structure and Adjustment Sets}
Figure~\ref{fig:dag} presents the directed acyclic graph (DAG) encoding our assumptions. Age, sex, education, occupation/employment, and noise exposure are treated as confounders of the perceived-stress$\rightarrow$outcome relationships and constitute the \textbf{minimal sufficient adjustment set} for the primary analysis. Depressed mood and sleep are treated as \emph{potential mediators} of the stress$\rightarrow$tinnitus path (and possibly colliders when conditioning opens non-causal paths); conditioning on them can constitute overadjustment. Accordingly:
\begin{itemize}[leftmargin=1.4em,itemsep=1pt]
\item \textbf{Primary model} adjusts for the minimal sufficient set only (age, sex, education, occupation/employment, noise).
\item \textbf{Extended model} additionally includes depressed mood and sleep, reported separately and interpreted as a mediator-adjusted (not confounder-adjusted) estimate.
\item \textbf{Mediation analysis} decomposes total, direct, and indirect (through depressed mood/sleep) associations under stated assumptions, as an exploratory addendum.
\item \textbf{Collider check} enumerates conditioning sets that could induce collider bias and excludes them from the primary specification.
\end{itemize}
Because restricting to employed participants can induce a healthy-worker effect (and a selection collider), the primary population is \emph{all eligible adults}, with an \emph{employed-restricted} analysis reported as a prespecified secondary contrast rather than the main result.

\subsection{Outcomes}
Endpoints are defined a priori and separated:
\begin{itemize}[leftmargin=1.4em,itemsep=1pt]
\item \textbf{Tinnitus presence} (any tinnitus in the past 12 months) and \textbf{bothersome tinnitus} (tinnitus causing annoyance/sleep or concentration difficulty) are modeled as \emph{distinct} endpoints, since their determinants differ.
\item \textbf{Hearing loss} is quantified by the speech-frequency pure-tone average (PTA$_{0.5,1,2,4}$) for \emph{both the better ear and the worse ear}, plus a \textbf{per-ear} (ear-level, mixed-effects) analysis, because a better-ear summary can mask the unilateral/asymmetric presentations most relevant to SSNHL. The primary threshold is PTA $>25$~dB HL (the WHO grade-1 cutoff \citep{who2021world}); $>20$~dB HL and high-frequency PTA$_{3,4,6}$ are prespecified sensitivity definitions, with the 25~dB HL choice justified as the established clinical cutpoint.
\item \textbf{Asymmetric hearing loss} is defined as an interaural difference $\geq$15~dB at $\geq$2 adjacent frequencies, consistent with asymmetry criteria used in sudden/asymmetric hearing-loss guidance \citep{chandrasekhar2019ssnhl}.
\item \textbf{Secondary:} self-reported hearing difficulty, hearing-aid use, and quality-of-life correlates.
\end{itemize}
\textbf{Noise exposure} is coded, where cycles permit, not merely as a binary but with duration, intensity/source (occupational, recreational, firearm), and hearing-protection use, so residual confounding by noise is minimized.

\subsection{Statistical Analysis}
All public-data analyses use survey-aware methods (Taylor-linearization variance estimation honoring weights, strata, and PSUs). Descriptive analyses estimate weighted prevalence with 95\% confidence intervals. Association analyses use weighted logistic, ordinal, or linear regression by outcome type. The \textbf{primary specification adjusts for the minimal sufficient set}; the extended (mediator-adjusted) model is secondary. Results are reported as associations (odds ratios or mean differences) with 95\% CIs, explicitly not as causal effects. Multiplicity is handled by prespecifying primary hypotheses (Introduction) and controlling the false-discovery rate across secondary contrasts.

\paragraph{Missing data.} Missingness is described by variable and mechanism. Primary analyses use multiple imputation by chained equations under a missing-at-random assumption (design variables in the imputation model), pooled by Rubin's rules; complete-case analysis is a sensitivity check. In external validation, records missing a common predictor are handled by the same prespecified rule in both surveys.

\paragraph{Sample size and precision.} Sample sizes are fixed, so we frame precision rather than post-hoc power: with several thousand audiometry participants per window, weighted prevalence is estimable to within a few percentage points and moderate associations (OR~$\approx$~1.3--1.5) are detectable at conventional error rates. Exact expected analytic $N$ per window is fixed in the repository before analysis.

\subsection{Research-Only Prediction Models and Intended Use}
Two targets---tinnitus presence and audiometric hearing loss---are developed as \emph{research risk-stratification models}, not screening, diagnostic, or triage devices. Their intended use, target user, prediction time, inputs, decision, and the costs of false positives/negatives and contraindicated uses are specified in Table~\ref{tab:intended}. We first fit transparent survey-weighted logistic baselines with prespecified predictors, then gradient-boosted trees as a comparator; the simpler model is preferred unless the complex model yields a clinically meaningful, calibrated improvement. Development uses KNHANES with repeated stratified cross-validation and optimism correction. All targets, predictors, and metrics are frozen before external validation (Table~\ref{tab:metrics}). SHAP values are reported for interpretability, with the explicit caveat that SHAP does not guarantee causal or even stable attributions when predictors are correlated; importance is interpreted descriptively, not causally.

\subsection{External Validation Across Survey Designs}
Transporting a KNHANES model to NHANES confronts differences in weighting, population prevalence, age range, category coding, and health-system context. We therefore prespecify:
\begin{itemize}[leftmargin=1.4em,itemsep=1pt]
\item \textbf{Unweighted vs.\ weighted performance} reported separately---unweighted for individual-level prediction, survey-weighted for population-level performance.
\item \textbf{Calibration-in-the-large and calibration slope reported separately}, before and after recalibration (intercept-only, then intercept+slope update of the frozen model).
\item \textbf{Common-support analysis} restricting comparison to the overlapping covariate region of the two surveys.
\item \textbf{Case-mix standardization} to distinguish genuine model transportability from case-mix differences in prevalence/age.
\item \textbf{A reduced common model} (predictors available and comparable in both surveys) versus \textbf{country-specific extended models}, so that limited measurement invariance does not masquerade as poor transportability.
\end{itemize}

\paragraph{Fairness and subgroup evaluation.} Discrimination and calibration are reported within prespecified subgroups: sex, age band, employment status, noise-exposure status, unilateral/asymmetric hearing, \emph{socioeconomic status} (education/income), and \emph{race/ethnicity} (NHANES, where collected). Subgroup calibration and error gaps are first-class results, not afterthoughts.

\subsection{Sensitivity Analyses}
Prespecified sensitivity analyses test alternative hearing-loss definitions (25 vs.\ 20~dB HL; high-frequency PTA; worse-ear and per-ear), age restrictions, exclusion of conductive patterns when tympanometry permits, alternative stress operationalizations, mediator- vs.\ confounder-treatment of depressed mood/sleep, all-adult vs.\ employed-restricted populations, and stratification by noise exposure, sex, and asymmetry.

\subsection{Ethics and Governance}
Stages~1--2 use publicly available, deidentified survey data and do not constitute human-subjects research requiring new consent; KNHANES and NHANES were approved by their respective review boards, and their use follows each provider's data-use terms. The Stage-3 clinical extension will obtain IRB approval at Ajou University Hospital, with a data-management plan covering deidentification, access control, and secondary-use consent. No identifiable patient data are placed in the public repository; only derived, deidentified, or synthetic artifacts are shared.
